
\subsection{Baseline BB84 Experiments} \label{sec:graph}

% Ideal, Noise-Free BB84

The BB84 protocol starts with the sender (Alice) encoding bits in a quantum state (qubits) and transmitting them over a (necessarily insecure) quantum channel. As described above, a physical realization of this is using vertical-horizontal or circular polarizations as a choice of basis each of which each have two orthogonal states in that basis. The receiver (Bob) must choose which basis to measure the incoming qubits without any knowledge of Alice's choice of basis. Bob has a 50\% chance of picking the correct basis, assuming the information is transmitted over a noiseless channel. After every message is sent, Alice and Bob publicly compare their choice of basis chosen for each message, and retain the bits where they both chose the same basis. If the public channel were secure and noiseless, then Alice and Bob would have a shared secret key that can be used for one-time pad encryption.

Figure \ref{fig:ideal} below shows an implementation of the BB84 protocol as a quantum circuit. The circuit starts with a quantum state in the Z-basis ($\left\{\ket{0}, \ket{1}\right\}$), that is then acted upon by an identity gate if Alice chooses to transmit in the z-basis, or is acted upon by a Pauli X gate if Alice chooses to transmit in the x-basis. The meter represents Bob's readout of the qubit into classical information.
\begin{figure}[ht]
    \centering
    \[
    \Qcircuit{
        \lstick{\text{Alice}} & \lstick{\ket{0/1}} & \gate{I/X} & \meter & \cw & \rstick{\text{Bob}}
    }
    \]
    \caption{Communication between Alice and Bob with no eavesdropper.}
    \label{fig:ideal}
\end{figure}


We implement three baseline scenarios of increasing complexity.
Our first baseline corresponds to the ideal case, in which neither noise nor eavesdropping is present in the quantum channel between Alice and Bob. In this setting, Alice’s state preparation, the quantum channel, and Bob’s measurements are assumed to be perfectly noiseless and free from external interference. This case is implemented only in local simulation, as current IBM quantum hardware is not fully fault-tolerant. Under these assumptions, Alice’s and Bob’s sifted keys are expected to be identical; equivalently, the quantum bit error rate (QBER) should be zero.

The second baseline is also a local simulation without Eve’s presence, but incorporates synthetic noise into the circuit via the simulate\_BB84 helper function. Specifically, we use a single-qubit depolarizing noise model with error probability 0.05. This model captures the effect of noise on Alice’s state preparation and Bob’s measurements, both of which involve single-qubit operations.

The final baseline executes the BB84 simulation on IBM’s quantum platforms. Since the specific backend is assigned automatically, we do not control which device is used. The goal in this case is to assess the level of noise in real hardware by comparing its performance against the ideal baseline.


